\documentclass[UTF8,11pt]{ctexart}
\usepackage[a4paper,margin=24mm]{geometry}
\usepackage{amsmath,amssymb,amsthm,mathtools,mathrsfs}
\usepackage[all]{xy}
\usepackage{xurl,hyperref,enumitem}
\usepackage{tikz}
\usetikzlibrary{arrows.meta,cd,decorations.markings,calc,patterns}
\hypersetup{hidelinks,breaklinks=true}
\setlength{\emergencystretch}{3em}
\allowdisplaybreaks
\newtheorem*{theorem}{Theorem}
\newtheorem*{thm}{Theorem}
\newtheorem*{lemma}{Lemma}
\newtheorem*{proposition}{Proposition}
\newtheorem*{prop}{Proposition}
\newtheorem*{corollary}{Corollary}
\newtheorem*{cor}{Corollary}
\newtheorem*{claim}{Claim}
\theoremstyle{definition}
\newtheorem*{definition}{Definition}
\newtheorem*{defn}{Definition}
\newtheorem*{remark}{Remark}
\newtheorem*{example}{Example}
\newcommand{\R}{\mathbb R}
\newcommand{\C}{\mathbb C}
\newcommand{\Z}{\mathbb Z}
\newcommand{\Q}{\mathbb Q}
\newcommand{\N}{\mathbb N}
\newcommand{\F}{\mathbb F}
\newcommand{\D}{\mathbb D}
\newcommand{\bB}{\mathbb B}
\newcommand{\bC}{\mathbb C}
\newcommand{\bR}{\mathbb R}
\newcommand{\bZ}{\mathbb Z}
\newcommand{\bN}{\mathbb N}
\newcommand{\bQ}{\mathbb Q}
\newcommand{\bD}{\mathbb D}
\newcommand{\bF}{\mathbb F}
\newcommand{\CP}{\mathbb{CP}}
\newcommand{\RP}{\mathbb{RP}}
\newcommand{\sO}{\mathscr O}
\newcommand{\sI}{\mathscr I}
\newcommand{\sF}{\mathscr F}
\newcommand{\sG}{\mathscr G}
\newcommand{\sH}{\mathscr H}
\newcommand{\sR}{\mathscr R}
\newcommand{\sM}{\mathscr M}
\newcommand{\sV}{\mathscr V}
\newcommand{\sU}{\mathscr U}
\DeclareMathOperator{\Hom}{Hom}
\DeclareMathOperator{\Ext}{Ext}
\DeclareMathOperator{\Tor}{Tor}
\DeclareMathOperator{\Spec}{Spec}
\DeclareMathOperator{\Proj}{Proj}
\DeclareMathOperator{\supp}{supp}
\DeclareMathOperator{\im}{im}
\DeclareMathOperator{\id}{id}
\DeclareMathOperator{\Tr}{Tr}
\DeclareMathOperator{\tr}{tr}
\DeclareMathOperator{\rank}{rank}
\DeclareMathOperator{\ord}{ord}
\DeclareMathOperator{\dist}{dist}
\DeclareMathOperator{\End}{End}
\DeclareMathOperator{\Aut}{Aut}
\DeclareMathOperator{\Gal}{Gal}
\DeclareMathOperator{\vol}{vol}
\DeclareMathOperator{\Cl}{Cl}
\DeclareMathOperator{\coker}{coker}
\DeclareMathOperator{\codim}{codim}
\DeclareMathOperator{\divergence}{div}
\DeclareMathOperator{\Ric}{Ric}
\DeclareMathOperator{\grad}{grad}
\DeclareMathOperator{\Hess}{Hess}
\newcommand{\abs}[1]{\left\lvert#1\right\rvert}
\newcommand{\sabs}[1]{\lvert#1\rvert}
\newcommand{\babs}[1]{\bigl\lvert#1\bigr\rvert}
\newcommand{\Babs}[1]{\Bigl\lvert#1\Bigr\rvert}
\newcommand{\bbabs}[1]{\biggl\lvert#1\biggr\rvert}
\newcommand{\BBabs}[1]{\Biggl\lvert#1\Biggr\rvert}
\newcommand{\norm}[1]{\left\lVert#1\right\rVert}
\newcommand{\snorm}[1]{\lVert#1\rVert}
\newcommand{\bnorm}[1]{\bigl\lVert#1\bigr\rVert}
\newcommand{\Bnorm}[1]{\Bigl\lVert#1\Bigr\rVert}
\newcommand{\bbnorm}[1]{\biggl\lVert#1\biggr\rVert}
\newcommand{\BBnorm}[1]{\Biggl\lVert#1\Biggr\rVert}
\newcommand{\widebar}[1]{\overline{#1}}
\newcommand{\nicefrac}[2]{\frac{#1}{#2}}
\newcommand{\linnprod}[2]{\langle#1,#2\rangle}
\newcommand{\myindex}[1]{#1}
\newcommand{\glsadd}[1]{}
\newcommand{\avoidbreak}{}
\newcommand{\sourceref}[1]{\textnormal{[原文：\nolinkurl{#1}]}}
\newcommand{\thmref}[1]{\sourceref{#1}}
\newcommand{\lemmaref}[1]{\sourceref{#1}}
\newcommand{\propref}[1]{\sourceref{#1}}
\newcommand{\corref}[1]{\sourceref{#1}}
\newcommand{\defnref}[1]{\sourceref{#1}}
\newcommand{\exampleref}[1]{\sourceref{#1}}
\newcommand{\exerciseref}[1]{\sourceref{#1}}
\newcommand{\figureref}[1]{\sourceref{#1}}
\newcommand{\sectionref}[1]{\sourceref{#1}}
\newcommand{\appendixref}[1]{\sourceref{#1}}
\newcommand{\stacksref}[2]{\href{https://stacks.math.columbia.edu/tag/#1}{#2 [#1]}}


\newcommand{\E}{\mathbb E}
\newcommand{\Prob}{\mathbb P}
\newcommand{\Reff}{\mathcal R_{\mathrm{eff}}}
\newcommand{\eps}{\varepsilon}
\DeclareMathOperator{\diam}{diam}
\DeclareMathOperator{\Var}{Var}
\DeclareMathOperator{\Crit}{Crit}
\newcommand{\dd}{\,\mathrm d}


\begin{document}
\section*{075\quad 拟随机有限群的卷积混合不等式}
\subsection*{来源与整理范围}
Terence Tao，\emph{254B, Notes 3: Quasirandom groups, expansion, and Selberg\textquotesingle s 3/16 theorem}，16 December 2011。Section 1，Proposition 3完整证明。
\par\url{https://terrytao.wordpress.com/2011/12/16/254b-notes-3-quasirandom-groups-expansion-and-selbergs-316-theorem/}
\par\textbf{恢复方式（整理者说明）：}从先前\nolinkurl{make074_075.py}执行记录恢复作者正文，重建独立排版；2026-09-28。
\subsection*{作者命题与人类证明}

\paragraph{Definition and convolution convention.}
A finite group $G$ is $D$-quasirandom if every non-trivial irreducible unitary representation of $G$ has dimension at least $D$. Given two functions $f,g\in\ell^2(G)$, we can define the convolution $f*g\in\ell^2(G)$ by the formula
\[f*g(x):=\sum_{y\in G}f(y)g(y^{-1}x)=\sum_{y\in G}f(xy^{-1})g(y).\]
\begin{proposition}[3, Mixing inequality]
Let $G$ be a finite $D$-quasirandom group, and let $f,g\in\ell^2(G)$. If at least one of $f,g$ has mean zero, then
\[\|f*g\|_{\ell^2(G)}\le D^{-1/2}|G|^{1/2}\|f\|_{\ell^2(G)}\|g\|_{\ell^2(G)}.\]
\end{proposition}
\begin{proof}
By subtracting a constant from $f$ or $g$, we may assume that $f$ or $g$ both have mean zero. Observe that $f*g$ (being a superposition of right-translates of $f$) also has mean zero. Thus, we see that we may define an operator $T_g:\ell^2(G)_0\to\ell^2(G)_0$ by setting $T_gf:=f*g$. It thus suffices to show that the operator norm of $T_g$ is at most $D^{-1/2}|G|^{1/2}\|g\|_{\ell^2(G)}$.

Fix $g$. We can view $T_g$ as a $|G|-1\times|G|-1$ matrix. We apply the singular value decomposition to this matrix to obtain singular values
\[\sigma_1\ge\ldots\ge\sigma_{|G|-1}\ge0\]
of $T_g$, together with associated singular vectors. The operator norm of $T_g$ is the largest singular value $\sigma_1$. The operator $T_g^*T_g$ is then a self-adjoint operator (or matrix) with eigenvalues $\sigma_1^2,\ldots,\sigma_{|G|-1}^2$. In particular, we have
\[\operatorname{tr}T_g^*T_g=\sigma_1^2+\ldots+\sigma_{|G|-1}^2.\]
Now, a short computation shows that $T_g^*T_gf=f*g*\tilde g$, where $\tilde g(x):=\overline{g(x^{-1})}$, and (by embedding $\ell^2(G)_0$ in $\ell^2(G)$, and noting that $T_g^*T_g$ annihilates constants) the trace can computed as
\[\operatorname{tr}T_g^*T_g=|G|g*\tilde g(0)=|G|\|g\|_{\ell^2(G)}^2.\]
Thus, if $V$ is the eigenspace of $T_g^*T_g$ corresponding to the eigenvalue $\sigma_1^2$ (so that the dimension of $V$ is the multiplicity of $\sigma_1$), we have
\[\dim(V)\sigma_1^2\le|G|\|g\|_{\ell^2(G)}^2.\]
Now observe that if $\tau:G\to U(\ell^2(G)_0)$ is the left-regular representation (restricted to $\ell^2(G)_0$) then
\[T_g^*T_g\tau(h)f=\tau(h)T_g^*T_gf\]
for any $f\in\ell^2(G)_0$ and $h\in G$ (this is a special case of the associativity of convolution). In particular, we see that $V$ is invariant under $\tau$. Since $\tau$ has no non-trivial invariant vectors in $\ell^2(G)_0$, we conclude from quasirandomness that $V$ has dimension at least $D$, and the claim follows.
\end{proof}

\subsection*{整理说明与先修范围（非作者正文）}
按先前生成脚本恢复。首段定义和卷积约定据原文Definition 1与Section 1摘编；该段为编辑合并，主命题及证明保持来源文字。$\ell^2$用计数测度，$\ell^2(G)_0$为零均值子空间。原文首句``f or g both与群单位元写作0的字样保留；此处0指卷积在群单位元的取值。只计卷积混合主命题，不把后面的乘积集合练习并入题面。标准先修是奇异值分解、迹和有限群酉表示；核心是识别最大特征空间的群不变性，以表示维数控制奇异值重数。
\subsection*{可修改的简短标注（非作者正文）}
$c$：有限群函数卷积的定量混合。

$a$：卷积算子；左正则表示；奇异值；迹；不变特征空间；不可约表示维数。

\texttt{cross\_theory\_status = yes}。将卷积误差转化为算子范数，利用群表示限制最大奇异值的重数，再用迹估计返回混合界。位置：Proposition 3。

全部标注为非穷尽、可修改初稿。
\subsection*{来源许可与编辑归属}
作者About页Copyright etc.允许复制合理部分（如单篇文章）并注明来源。本题依该许可摘录：\url{https://terrytao.wordpress.com/about/}。
ChatGPT 加入题号、中文来源说明、整理范围和初稿标注；编辑说明与人类证明分列。
\end{document}
